%%%%%
%
\section{Longueur de Monin-Obukhov (formule exacte)}

\itemstyle{Op\'erateur~:} {\sf Mol\_exact}
\index{Mol\_exact}
\newline

%%%%%OBJET%%%%
\itemstyle{Objet~:}
Longueur de Monin-Obukhov calcul\'ee par la formule exacte � partir des composantes du vent, des temp\'eratures (sonique, de l'air) et du rapport de m\'elange. La temp�rature sonique est assimil�e � la temp�rature potentielle virtuelle $\theta_{v}$ ; cette approximation est l�gitime car la temp�rature dite sonique est tr�s proche de $T_{v}$ elle m�me tr�s proche de $\theta_{v}=T_{v}.(\frac{P_{0}}{P})^{0.286}$ quand la pression au point de mesure est peu diff�rente de $P_{0}$ (avec $P_{0}=1000$ hPa). LMO est la hauteur (entre -200 et +200 m) � laquelle la production thermique de turbulence l'emporte sur la production m�canique :
\begin{center}
L>0 : stabilit� ; L=0 : convection libre ; L<0 : instabilit� 
\end{center}

%%%%%ARGUMENTS%%%%
\bigskip
%%%%%ENTREES%%%%
\begin{tabarg}
\multicolumn{3}{l}{\titretab {Entr\'ees}}\\
\hline
$U_{f}$ & Vecteur, \texttt{float} & 1�re composante horizontale du vent filtr�e type FRC300 (\textit{m/s})\\
$V_{f}$ & Vecteur, \texttt{float} & 2�me composante horizontale du vent (\textit{m/s})\\
$W_{f}$ & Vecteur, \texttt{float} & 3�me composante verticale du vent (\textit{m/s})\\
$T_{sonique}$ & Vecteur, \texttt{float} & Temp\'erature sonique non filtr�e (\textit{C} ou \textit{K})\\
$T_{air}$ & Vecteur, \texttt{float} & Temp\'erature de l'air d\'eduite du Sonique et du Licor non filtr�e (\textit{C} ou \textit{K})\\
$r$ & Vecteur, \texttt{float} & Rapport de m\'elange en vapeur d'eau non filtr�(\textit{g/kg})\\
$Fe$ & Coeff., \texttt{float} & Fr\'equence d'\'echantillonnage des donn\'ees (\textit{/s})\\
$t_{int}$ & Coeff., \texttt{float} & Temps d'int\'egration pour le calcul (\textit{s})\\
$tau$ & Coeff., \texttt{float} & Constante de temps du filtrage (g�n�ralement entre 200 et 600 sec)(\textit{s})\\
$rap_{mnq}$ & Coeff., \texttt{float} & Souvent 0.5 si plus de la moiti\'e des points (\textit{any})\\
\end{tabarg}

\medskip
%%%%%SORTIES%%%%
\begin{tabarg}
\multicolumn{3}{l}{\titretab {Sorties}}\\
\hline
$L_{mon}$ & Vecteur, \texttt{float} & Longueur de Monin Obukhov (\textit{m})\\
\end{tabarg}


%%%%%FORMULE%%%%
\itemstyle{Formule}

%%Mol_exact

Le r\'esultat $L_{mon}$ est obtenu � partir du calcul des covariances.

\begin{align*}
u_* &= ({\overline{w'v'}}^{2}+{\overline{w'v'}}^{2})^{\frac{1}{4}}\\
\intertext{On en d\'eduit :}
L_{mon} &= \frac{-u_*^3.\overline{T_{v}}}{k.g.\overline{w'T_{vf}'}}                   \end{align*}

avec \newline 
        $k$ : constante de Kolmogorov choisie �gale � 0.4\newline
        $g$ : constante de gravit� �gale � 9.81 $ms^{-2}$




%%%%%AUTEUR%%%%
\itemstyle{Auteur~:}
CNRM/GMEI/TRAMM

%%%%%REFERENCE%%%%
\itemstyle{R\'ef\'erences~:}
